\begin{center}
{\Huge \scshape Qiong Liu, Ph.D.} \\ \vspace{4pt}
\small
\faPhone\ 475-280-1364 \quad
\href{mailto:qiongliu.work@gmail.com}{qiongliu.work@gmail.com} \quad
\href{https://www.linkedin.com/in/qiong-l-529b23196}{LinkedIn} \quad
\href{https://scholar.google.com/citations?user=SvnpHCkAAAAJ&hl=en}{Google Scholar}

Authorized to work in the U.S. without sponsorship
\end{center}

\section{Professional Summary}
\textbf{Machine Learning Engineer} with experience building end-to-end ML pipelines for image quality improvement, signal restoration, and robust modeling under noisy and incomplete data conditions. Strong background in \textbf{deep learning, self-supervised learning, evaluation design, and real-world data analysis}. Proven track record of translating technically complex problems into deployable algorithmic solutions across research and product-oriented settings.

\section{Core Expertise}
\textbf{Applied ML:} Deep learning, self-supervised learning, multimodal learning, diffusion models, transformers \\
\textbf{Image / Signal Quality:} Denoising, restoration, motion artifact reduction, quality evaluation, robustness analysis \\
\textbf{Modeling \& Evaluation:} Spatiotemporal modeling, representation learning, uncertainty modeling, quantitative validation \\
\textbf{Programming:} Python, PyTorch, NumPy, SciPy, MATLAB

\section{Experience}
\textbf{Canon Medical Research USA} \hfill Apr 2025 -- Present \\
\textit{AI Scientist (Reconstruction), Vernon Hills, IL}
\begin{itemize}[leftmargin=*]
\item Built machine learning pipelines for \textbf{image quality improvement} and motion-related degradation correction in real-world imaging data
\item Developed robust algorithms for noisy, low-signal, and incomplete data settings, improving downstream consistency and usability
\item Designed quantitative evaluation workflows to assess model performance, generalization, and reliability across data conditions
\item Worked cross-functionally to move algorithm development toward production-oriented deployment
\end{itemize}

\textbf{Canon Medical Research USA} \hfill Jun 2023 -- May 2024 \\
\textit{Research Scientist Intern, Vernon Hills, IL}
\begin{itemize}[leftmargin=*]
\item Developed deep learning models for \textbf{signal restoration and denoising}, preserving critical structures under challenging acquisition conditions
\item Designed validation frameworks for robustness, consistency, and quality assessment across heterogeneous datasets
\item Communicated technical results to support algorithm design and product decisions
\end{itemize}

\textbf{United Imaging Healthcare} \hfill Jul 2019 -- Aug 2019 \\
\textit{Student Intern, Wuhan, China}
\begin{itemize}[leftmargin=*]
\item Supported development and validation of engineering systems for clinical applications
\end{itemize}

\section{Patents}
\textbf{Edge-Guided Deformation Field Fusion for Motion Correction} \hfill 2026 \\
Patent Pending, \textit{Lead Inventor}
\begin{itemize}[leftmargin=*]
\item Developed a spatially adaptive algorithm to improve robustness and alignment quality in complex image regions
\end{itemize}

\textbf{Automatic Data-Driven Multi-Signal Decomposition} \hfill 2025 \\
Patent Pending, \textit{Lead Inventor}
\begin{itemize}[leftmargin=*]
\item Developed an ML-based method to separate overlapping temporal signals from unlabeled dynamic data
\end{itemize}

\section{Selected Machine Learning Projects}

\textbf{Robust Motion-Aware Image Modeling}
\begin{itemize}[leftmargin=*]
\item Developed ML methods for motion-related image degradation correction under noisy, low-signal, and heterogeneous real-world conditions
\item Proposed spatially adaptive regularization and structured spatial guidance to improve model robustness beyond conventional globally constrained approaches
\item Improved consistency and reliability of downstream image analysis in challenging temporal data
\end{itemize}

\textbf{Self-Supervised Temporal Denoising}
\begin{itemize}[leftmargin=*]
\item Built self-supervised denoising models for temporal imaging data with explicit consistency constraints across frames
\item Incorporated domain-informed priors into training to improve robustness, generalization, and preservation of meaningful temporal structure
\item Demonstrated that structured constraints can improve model performance in noisy and limited-data regimes
\end{itemize}

\textbf{Interpretable Deep Learning for Restoration}
\begin{itemize}[leftmargin=*]
\item Developed CNN-based and deep image prior models for denoising and reconstruction across heterogeneous data conditions
\item Studied the effects of training data composition and optimization behavior on restoration quality and model reliability
\item Analyzed model behavior during training to better understand how deep networks learn signal and noise structure
\end{itemize}

\section{Education}
\textbf{Yale University} \hfill Aug 2020 -- May 2025 \\
Ph.D. in Biomedical Engineering \\
Thesis: Improving Image Quality and Quantification Accuracy for Static and Dynamic PET

\textbf{Huazhong University of Science and Technology} \hfill Sep 2016 -- Jun 2020 \\
B.S. in Biomedical Engineering, Minor in Computer Science \\
GPA: 3.94/4.00

\section{Selected Publications}
\begin{itemize}[leftmargin=*]
\item Liu Q., et al., \textit{Patlak-guided self-supervised learning for dynamic PET denoising}, \textbf{IEEE Transactions on Radiation and Plasma Medical Sciences}, 2026
\item Liu Q., et al., \textit{Population-based deep image prior for dynamic PET denoising: a data-driven approach to improve parametric quantification}, \textbf{Medical Image Analysis}, 2024
\item Liu Q., et al., \textit{A personalized deep learning denoising strategy for low-count PET images}, \textbf{Physics in Medicine \& Biology}, 2022
\end{itemize}

\section{Awards}
\textbf{Society of Nuclear Medicine and Molecular Imaging (SNMMI)}
\begin{itemize}[leftmargin=*]
\item Young Investigator Award (1st Place), 2025
\item Poster Award (2nd Place), 2024
\item Young Investigator Award (2nd Place), 2023
\end{itemize}

\textbf{Mitacs}
\begin{itemize}[leftmargin=*]
\item Globalink Research Award, 2019
\end{itemize}